
Un reloj de péndulo metálico adelanta $\tau_1$=5 s por día a una temperatura \textit{t}$_1$=+15 $^o$C y atrasa $\tau_2$ = 10 s por día a una temperatura \textit{t}$_2$=30 $^o$C. Encontrar el coeficiente $\alpha$ de expansión térmica del metal del péndulo, considerando que el período de oscilación de éste $T = 2\pi \sqrt{l/g}$ (\textit{l} es la longitud del péndulo; \textit{g}, la aceleración de la caída libre).